\section{Introduction}
\label{sec:intro}

\todo{Motivation paragraph: why cross-dataset emotion matters. ~5 sentences.}

Emotion classification has been studied extensively on individual
corpora, but \emph{cross-dataset} performance under realistic
deployment — where the training distribution differs from the
inference distribution — has received less systematic attention.
Three obstacles compound the difficulty: (i) label-space
incompatibility across emotion datasets, which forces either
projection to a shared schema or per-corpus reannotation
\citep{bostan2018analysis}; (ii) severe and often per-corpus class
imbalance, including outright missing classes (e.g.\ ISEAR contains no
\emph{surprise} examples); and (iii) register shift between corpora
collected from different sources (Reddit comments, self-report
narratives, empathy essays).

\todo{Problem statement paragraph. Why prior solutions are insufficient. ~5 sentences.}

The standard answer in the cross-domain literature has been
domain-adversarial training and its extensions \citep{ganin2016dann,
long2018cdan}. These methods were validated on cross-domain sentiment
\citep{glorot2011sentiment} and multi-domain text classification
\citep{chen2018multinomial}, but their transfer to cross-dataset
emotion has not been systematically tested. Concurrently, focal loss
\citep{lin2017focal} has been applied to imbalanced classification
across image and text domains, including emotion-specific imbalance
\citep{suresh2021labelaware}, but not in combination with domain-adversarial
frameworks on cross-dataset emotion benchmarks.

\todo{Contributions paragraph + finding teasers. ~7 sentences.}

We make three contributions. First, we provide the first systematic
six-method comparison (CE source-only, mixed-CE, DANN, CDAN, DANN+Focal,
CDAN+Focal) on three Ekman-mapped emotion corpora under two
evaluation protocols (mixed training and leave-one-dataset-out),
with three to five seeds per cell and paired bootstrap significance
testing.
Second, we report a counter-intuitive empirical finding: the simplest
non-adversarial method — source-only training with focal loss
replacing cross-entropy — achieves the largest measured cross-dataset
gain (\focalwassagain{} F1 points on WASSA-21 LODO,
$p=\focalwassap$). Adding domain-adversarial alignment to focal
either preserves or reduces this gain, never amplifies it.
Third, we characterize two distinct failure modes for adversarial
methods in this regime: DANN exhibits training instability at moderate
$\lambda$ values (gradient-reversal scale), and CDAN, while stable,
enforces source-source alignment that does not extrapolate to the
held-out target. Two of three LODO targets are
method-invariant: no pair of methods differs significantly on
GoEmotions LODO or ISEAR LODO.

We argue that the dominant cross-dataset emotion shift is
\emph{class-distribution shift} (which class frequencies and presence
differ across corpora), not \emph{feature-distribution shift} (how
classes are represented), and that this mismatch explains why
adversarial methods underperform a class-imbalance-aware loss on the
same task.

\paragraph{Paper organization.}
Section~\ref{sec:related} surveys the three literature threads that
motivate our experimental matrix. Section~\ref{sec:methods}
describes datasets, label mapping, six methods, and two protocols.
Section~\ref{sec:setup} details hyperparameters and evaluation.
Section~\ref{sec:results} presents the main results, bootstrap
significance tests, and the DeBERTa-large robustness validation.
Section~\ref{sec:discussion} analyzes failure modes and the
class-imbalance hypothesis. The Limitations section discusses
single-seed cells, target test-set size, and scope.
